\chapter{Réfraction de la lumière}\label{ch:g10:refraction}

Plongez une paille dans un verre d'eau~: elle paraît cassée à la surface.
Les piscines sont toujours plus profondes qu'elles n'en ont l'air, les
diamants jettent un feu coloré depuis une pierre incolore, et la phrase
que vous lisez a probablement traversé un océan dans un fil de verre
plus fin qu'un cheveu. Une seule loi explique les quatre --- une égalité
entre deux sinus, trouvée par Snell en 1621 après quinze siècles de
physiciens contemplant des tables d'angles sans voir le motif.

\section{Réflexion~: un rappel}

Dans un milieu transparent homogène, la lumière se propage en ligne
droite --- les rayons du~\cref{ch:g10:light-spectra}. À la frontière
entre deux milieux, une partie de la lumière rebondit
(\emph{réflexion}) et une partie passe (\emph{\omterm{def:g10:refraction:refraction}{réfraction}}). Les deux
obéissent à des lois géométriques exactes, énoncées avec une
convention~: les angles se mesurent depuis la
\emph{\omterm{def:g10:refraction:angles}{normale}}, jamais depuis la surface.

\begin{definition}[Normale et angle d'incidence]\label{def:g10:refraction:angles}
Au point où un rayon rencontre une surface, la
\emph{normale}\index{normale} est la droite perpendiculaire à la
surface en ce point. L'\emph{angle d'incidence}\index{angle
d'incidence} $i_1$ est l'angle entre le rayon incident et la
normale~; les angles de réflexion et de \omterm{def:g10:refraction:refraction}{réfraction} se mesurent
depuis la normale de la même façon.
\end{definition}

\begin{proposition}[Loi de la réflexion]\label{prop:g10:refraction:reflection}
Le rayon réfléchi se trouve dans le plan contenant le rayon incident et
la \omterm{def:g10:refraction:angles}{normale}, de l'autre côté de la \omterm{def:g10:refraction:angles}{normale}, et l'angle de
réflexion égale l'\omterm{def:g10:refraction:angles}{angle d'incidence}~:
\[
i_1' = i_1 .
\]
\end{proposition}
\admitted

\begin{example}[Lire correctement les angles]\label{ex:g10:refraction:mirror}
Un faisceau laser frappe un miroir «~à $35^\circ$~» --- mesuré depuis la
surface du miroir, comme on décrit souvent les faisceaux. L'\omterm{def:g10:refraction:angles}{angle
d'incidence} est $90^\circ - 35^\circ = 55^\circ$, donc le rayon
réfléchi part aussi à $55^\circ$ de la \omterm{def:g10:refraction:angles}{normale}, et le faisceau
tourne de $180^\circ - 2 \times 55^\circ = 70^\circ$. Oublier la
convention de la \omterm{def:g10:refraction:angles}{normale} est l'erreur la plus courante de ce
chapitre~; elle coûte $90^\circ$ à chaque fois.
\end{example}

\section{Réfraction et indice de réfraction}

Tenez un crayon dans un verre à moitié rempli~: il paraît cassé à la
surface. La lumière de la partie immergée change de direction en
quittant l'eau, donc l'œil --- qui suppose des rayons droits ---
reconstruit la pointe quelque part où elle n'est pas.

\begin{definition}[Réfraction]\label{def:g10:refraction:refraction}
La \emph{réfraction}\index{réfraction} est le changement de direction de
la lumière lorsqu'elle traverse la frontière entre deux milieux
transparents. Le rayon dans le second milieu est le \emph{rayon
réfracté}\index{rayon réfracté}, et son angle depuis la
\omterm{def:g10:refraction:angles}{normale} est l'angle de réfraction $i_2$.
\end{definition}

La \omterm{def:g10:refraction:refraction}{réfraction} se produit parce que la lumière a des vitesses
différentes dans des milieux différents~; chaque matériau se
caractérise par le ralentissement qu'il impose à la lumière.

\begin{definition}[Indice de réfraction]\label{def:g10:refraction:index}
L'\emph{indice de réfraction}\index{indice de réfraction} d'un milieu
transparent est le rapport
\[
n = \frac{c}{v},
\]
où $c = \qty{3.00e8}{m/s}$ est la vitesse de la lumière dans le vide et
$v$ sa vitesse dans le milieu. Comme $v \leq c$, l'indice satisfait
$n \geq 1$~; il n'a pas d'\omterm{def:g10:orders-of-magnitude:unit}{unité}. Quelques valeurs utiles~:
\begin{center}
\begin{tabular}{lccccc}
milieu & air & eau & plexiglas & verre & diamant \\
$n$ & $1.0003$ & $1.33$ & $1.49$ & $1.50$ & $2.42$ \\
\end{tabular}
\end{center}
\end{definition}

\begin{example}[Vitesse de la lumière dans l'eau]\label{ex:g10:refraction:speed}
Dans l'eau, $v = \dfrac{c}{n} = \dfrac{\qty{3.00e8}{m/s}}{1.33}
= \qty{2.26e8}{m/s}$. Dans le diamant, la lumière rampe à
$\qty{3.00e8}{m/s}/2.42 = \qty{1.24e8}{m/s}$ --- moins de la moitié de
sa vitesse dans le vide. Pour l'air, $n = 1.0003$ change la vitesse de
$0.03\%$~: dans ce chapitre on traite l'air comme le vide et on prend
$n_{\text{air}} = 1.00$.
\end{example}

\begin{theorem}[Loi de Snell--Descartes de la réfraction]\label{thm:g10:refraction:snell}
\index{loi de Snell--Descartes}
Le \omterm{def:g10:refraction:refraction}{rayon réfracté} se trouve dans le plan du rayon incident et de
la \omterm{def:g10:refraction:angles}{normale}, de l'autre côté de la \omterm{def:g10:refraction:angles}{normale}, et les angles
d'incidence et de \omterm{def:g10:refraction:refraction}{réfraction} satisfont
\[
n_1 \sin i_1 = n_2 \sin i_2 ,
\]
où $n_1$ et $n_2$ sont les indices de réfraction des deux milieux. Le
trajet est réversible~: la lumière voyageant en sens inverse le long du
\omterm{def:g10:refraction:refraction}{rayon réfracté} sort le long du rayon incident.
\end{theorem}
\admitted

\begin{remark}\label{rem:g10:refraction:honest}
Les lois de la réflexion et de la \omterm{def:g10:refraction:refraction}{réfraction} sont énoncées ici comme
faits expérimentaux. Elles sont dérivées honnêtement dans le volume de
1\textsuperscript{re} année, où la lumière est traitée comme une
onde~: les deux lois --- et la formule $n = c/v$ elle-même ---
suivent alors du ralentissement de l'onde dans le milieu plus dense.
\end{remark}

\begin{omfigure}
\begin{tikzpicture}[scale=1.0, font=\small]
  \fill[blue!8] (-3.2,-2.1) rectangle (3.2,0);
  \draw[thick] (-3.2,0) -- (3.2,0);
  \draw[dashed] (0,-2.1) -- (0,2.3);
  \coordinate (O) at (0,0);
  \coordinate (I) at (-1.8,2.0);
  \coordinate (R) at (1.8,2.0);
  \coordinate (T) at (1.0,-1.73);
  \coordinate (N) at (0,2.3);
  \coordinate (Nd) at (0,-2.1);
  \draw[-Stealth, omDef, thick] (I) -- (O);
  \draw[-Stealth, omExo, thick] (O) -- (R);
  \draw[-Stealth, omThm, thick] (O) -- (T);
  \pic[draw, angle radius=9mm, "$i_1$", angle eccentricity=1.35] {angle = N--O--I};
  \pic[draw, angle radius=7mm, "$i_1'$", angle eccentricity=1.4] {angle = R--O--N};
  \pic[draw, angle radius=9mm, "$i_2$", angle eccentricity=1.35] {angle = Nd--O--T};
  \node[anchor=west] at (-3.1,0.45) {milieu 1 ($n_1$)};
  \node[anchor=west] at (-3.1,-0.55) {milieu 2 ($n_2$)};
\end{tikzpicture}

{\small Rayons incident (bleu), réfléchi (gris) et réfracté (rouge).
Tous les angles se mesurent depuis la \omterm{def:g10:refraction:angles}{normale} (en tirets)~; en
entrant dans le milieu plus dense ($n_2 > n_1$), le rayon se courbe
\emph{vers} la \omterm{def:g10:refraction:angles}{normale}.}
\end{omfigure}

\begin{example}[De l'air à l'eau]\label{ex:g10:refraction:airwater}
Un rayon frappe un étang calme à $i_1 = 30^\circ$. Avec $n_1 = 1.00$ et
$n_2 = 1.33$~:
\[
\sin i_2 = \frac{n_1 \sin i_1}{n_2} = \frac{\sin 30^\circ}{1.33}
= \frac{0.500}{1.33} = 0.376,
\qquad
i_2 = 22.1^\circ .
\]
Le rayon se courbe vers la \omterm{def:g10:refraction:angles}{normale}, comme toujours lorsqu'il entre
dans un milieu plus lent ($n_2 > n_1$). En quittant l'eau, le même
calcul se fait à l'envers et le rayon s'écarte \emph{de} la
\omterm{def:g10:refraction:angles}{normale}.
\end{example}

\begin{method}[Résoudre un problème de réfraction]\label{met:g10:refraction:method}
\begin{enumerate}
  \item Dessiner la surface, la \omterm{def:g10:refraction:angles}{normale} et le rayon~; identifier
        $n_1$ (milieu du rayon incident) et $n_2$.
  \item Vérifier que les angles donnés sont mesurés depuis la
        \omterm{def:g10:refraction:angles}{normale}~; convertir s'ils le sont depuis la surface.
  \item Écrire $n_1 \sin i_1 = n_2 \sin i_2$ et isoler le sinus
        inconnu~; obtenir l'angle avec le sinus inverse de la
        calculatrice, comme dans le volume de mathématiques.
  \item Contrôle de vraisemblance~: le rayon se courbe vers la
        \omterm{def:g10:refraction:angles}{normale} si $n_2 > n_1$, s'en écarte si $n_2 < n_1$, et
        $\sin i_2$ doit sortir entre $0$ et $1$.
\end{enumerate}
\end{method}

L'expérience derrière la loi tient sur un bureau~: un demi-disque de
plexiglas, un rapporteur, un laser. Les paires d'angles brutes
$(i_1, i_2)$ paraissent irrégulières --- Ptolémée les a tabulées vers
l'an~150 et a conclu à tort que $i_2$ était proportionnel à $i_1$.
Tracer le sinus contre le sinus, et les données s'alignent en une
droite passant par l'origine de pente $n_1/n_2$~: cette droite \emph{est}
la \omterm{thm:g10:refraction:snell}{loi de Snell--Descartes}.

\begin{omfigure}
\begin{tikzpicture}
\begin{axis}[omaxis, width=7.2cm, height=5.4cm,
  xlabel={$\sin i_1$}, ylabel={$\sin i_2$},
  xmin=0, xmax=1.08, ymin=0, ymax=0.88,
  xtick={0,0.2,0.4,0.6,0.8,1.0}, ytick={0.2,0.4,0.6,0.8}]
  \addplot[omThm, thick, domain=0:1.02] {x/1.33};
  \addplot[only marks, mark=*, mark size=1.6pt, omDef] coordinates
    {(0.174,0.131) (0.342,0.257) (0.500,0.376) (0.643,0.483)
     (0.766,0.576) (0.866,0.651) (0.940,0.707)};
\end{axis}
\end{tikzpicture}

{\small Mesures air--eau pour $i_1 = 10^\circ$ à $70^\circ$~: sinus contre
sinus, les données s'alignent sur une droite passant par l'origine de
pente $1/1.33$.}
\end{omfigure}

\begin{proposition}[Profondeur apparente]\label{prop:g10:refraction:depth}
\index{profondeur apparente}
Un objet à la profondeur $d$ sous l'eau, vu d'en haut sous de petits
angles, paraît se trouver à la \emph{\omterm{prop:g10:refraction:depth}{profondeur apparente}}
\[
d' = \frac{d}{n} ,
\]
où $n$ est l'indice de l'eau. Une piscine de profondeur \qty{1.20}{m}
ne paraît profonde que de $1.20/1.33 = \qty{0.90}{m}$ --- cause
classique de plongeons imprudents.
\end{proposition}
\admitted

\begin{remark}\label{rem:g10:refraction:depthhonest}
La formule suit de la \omterm{thm:g10:refraction:snell}{loi de Snell--Descartes} appliquée aux
rayons presque verticaux qui atteignent l'œil~; la dérivation est faite
dans le volume de 1\textsuperscript{re} année. Le sens de l'effet n'a
besoin d'aucune formule~: les rayons quittant l'eau s'écartent de la
\omterm{def:g10:refraction:angles}{normale}, donc l'œil les prolonge en arrière vers un point plus
haut que l'objet.
\end{remark}

\section{Dispersion~: une loi par couleur}

Un prisme transforme la lumière solaire blanche en arc-en-ciel, comme
les spectres du~\cref{ch:g10:light-spectra} l'ont montré. La
\omterm{def:g10:refraction:refraction}{réfraction} explique \emph{pourquoi}~: l'\omterm{def:g10:refraction:index}{indice de
réfraction} d'un milieu n'est pas tout à fait le même pour chaque
couleur.

\begin{definition}[Dispersion]\label{def:g10:refraction:dispersion}
Un milieu est \emph{dispersif}\index{dispersion} lorsque son \omterm{def:g10:refraction:index}{indice
de réfraction} dépend de la \omterm{def:g10:light-spectra:wavelength}{longueur d'onde} de la lumière. Dans le
verre et l'eau, $n$ est légèrement plus grand pour la lumière bleue
(\omterm{def:g10:light-spectra:wavelength}{courte longueur d'onde}) que pour la lumière rouge
(\omterm{def:g10:light-spectra:wavelength}{longue longueur d'onde}) --- donc le bleu se courbe davantage.
\end{definition}

\begin{example}[Deux couleurs, une surface]\label{ex:g10:refraction:twocolors}
Un faisceau de \omterm{def:g10:light-spectra:white}{lumière blanche} frappe du crown glass à
$i_1 = 60^\circ$. Pour ce verre $n_{\text{red}} = 1.51$ et
$n_{\text{blue}} = 1.53$. La \omterm{thm:g10:refraction:snell}{loi de Snell--Descartes}, une fois par
couleur~:
\[
\sin i_{\text{red}} = \frac{\sin 60^\circ}{1.51} = 0.574,
\quad i_{\text{red}} = 35.0^\circ ;
\qquad
\sin i_{\text{blue}} = \frac{\sin 60^\circ}{1.53} = 0.566,
\quad i_{\text{blue}} = 34.5^\circ .
\]
Une seule surface sépare la \omterm{def:g10:light-spectra:white}{lumière blanche} d'un demi-degré ---
invisible sur une vitre, mais un prisme réfracte deux fois dans le même
sens de rotation, élargit l'éventail, et à quelques mètres les couleurs
sont séparées de centimètres.
\end{example}

\begin{remark}\label{rem:g10:refraction:rainbow}
L'arc-en-ciel est le même calcul mené à l'intérieur d'une goutte de
pluie~: la lumière solaire se réfracte en entrant dans la goutte, se
réfléchit sur sa paroi arrière, et se réfracte à nouveau en sortant.
L'indice de l'eau va de $1.331$ (rouge) à $1.343$ (bleu), donc chaque
couleur sort sous son propre angle --- et aucun observateur ne voit le
même arc-en-ciel.
\end{remark}

\section{Réflexion totale interne}

La \omterm{thm:g10:refraction:snell}{loi de Snell--Descartes} contient un piège. En allant de
l'eau vers l'air, le \omterm{def:g10:refraction:refraction}{rayon réfracté} s'écarte de la
\omterm{def:g10:refraction:angles}{normale}~: $\sin i_2 = 1.33 \sin i_1 > \sin i_1$. Augmentez
$i_1$, et la formule exige bientôt $\sin i_2 > 1$, qu'aucun angle ne
peut livrer. La \omterm{def:g10:refraction:refraction}{réfraction} cesse alors purement et simplement.

\begin{definition}[Angle critique et réflexion totale interne]\label{def:g10:refraction:critical}
Soit de la lumière voyageant dans un milieu d'indice $n_1$ vers un
milieu d'indice plus petit $n_2 < n_1$. L'\emph{angle
critique}\index{angle critique} $i_c$ est l'\omterm{def:g10:refraction:angles}{angle d'incidence} pour
lequel le \omterm{def:g10:refraction:refraction}{rayon réfracté} rase la surface ($i_2 = 90^\circ$).
Pour $i_1 > i_c$ aucun \omterm{def:g10:refraction:refraction}{rayon réfracté} n'existe~: la surface
réfléchit toute la lumière dans le premier milieu, suivant la loi de la
réflexion. C'est la \emph{réflexion totale interne}\index{réflexion
totale interne}.
\end{definition}

\begin{proposition}[Formule de l'angle critique]\label{prop:g10:refraction:criticalformula}
Pour $n_1 > n_2$,
\[
\sin i_c = \frac{n_2}{n_1} .
\]
\end{proposition}

\begin{proof}
Poser $i_2 = 90^\circ$ dans la \omterm{thm:g10:refraction:snell}{loi de Snell--Descartes}~:
$n_1 \sin i_c = n_2 \sin 90^\circ = n_2$. Pour $i_1 > i_c$, la
\omterm{thm:g10:refraction:snell}{loi de Snell--Descartes} exigerait
$\sin i_2 = \frac{n_1}{n_2}\sin i_1 > 1$~: aucune direction réfractée
n'existe. (Que la lumière soit alors \emph{entièrement} réfléchie ---
non absorbée --- est un fait expérimental, confirmé par la théorie
ondulatoire du volume de 1\textsuperscript{re} année.)
\end{proof}

\begin{omfigure}
\begin{tikzpicture}[scale=0.92, font=\small]
  \fill[blue!8] (-3.6,-1.9) rectangle (3.6,0);
  \draw[thick] (-3.6,0) -- (3.6,0);
  \draw[dashed] (-2.2,-1.6) -- (-2.2,1.6);
  \draw[dashed] (0.3,-1.6) -- (0.3,1.6);
  \draw[dashed] (2.2,-1.6) -- (2.2,1.6);
  \draw[-Stealth, omDef, thick] (-2.95,-1.30) -- (-2.2,0);
  \draw[-Stealth, omThm, thick] (-2.2,0) -- (-1.20,1.12);
  \draw[-Stealth, omDef, thick] (-0.83,-0.99) -- (0.3,0);
  \draw[-Stealth, omThm, thick] (0.3,0) -- (1.75,0.06);
  \draw[-Stealth, omDef, thick] (0.90,-0.75) -- (2.2,0);
  \draw[-Stealth, omDef, thick] (2.2,0) -- (3.50,-0.75);
  \node at (-2.2,-1.85) {$i_1 < i_c$};
  \node at (0.3,-1.85) {$i_1 = i_c$};
  \node at (2.2,-1.85) {$i_1 > i_c$};
\end{tikzpicture}

{\small De l'eau vers l'air, à incidence croissante~: \omterm{def:g10:refraction:refraction}{réfraction} en
s'écartant de la \omterm{def:g10:refraction:angles}{normale}, puis un \omterm{def:g10:refraction:refraction}{rayon réfracté} rasant la
surface à l'\omterm{def:g10:refraction:critical}{angle critique}, puis \omterm{def:g10:refraction:critical}{réflexion totale
interne} --- la surface est devenue un miroir parfait.}
\end{omfigure}

\begin{example}[Trois angles critiques]\label{ex:g10:refraction:criticalvalues}
Vers l'air~: pour l'eau, $\sin i_c = 1/1.33 = 0.752$, donc
$i_c = 48.8^\circ$~; pour le verre, $\sin i_c = 1/1.50$, donc
$i_c = 41.8^\circ$~; pour le diamant, $\sin i_c = 1/2.42 = 0.413$,
donc $i_c = 24.4^\circ$. Plus l'indice est élevé, plus le cône
d'échappement est étroit~: dans le diamant, tout rayon à plus de
$24.4^\circ$ de la \omterm{def:g10:refraction:angles}{normale} d'une facette est piégé --- le secret
de son éclat, disséqué dans le~\cref{pb:g10:refraction:1}.
\end{example}

\section{Fibres optiques}

La \omterm{def:g10:refraction:critical}{réflexion totale interne} transforme un fil de verre en
tuyau pour la lumière~: un rayon entrant presque parallèle à l'axe
frappe la paroi bien au-delà de l'\omterm{def:g10:refraction:critical}{angle critique}, se réfléchit sans
perte, et le fait des millions de fois par kilomètre sans fuir.

\begin{definition}[Fibre optique]\label{def:g10:refraction:fiber}
Une \emph{fibre optique}\index{fibre optique} est un fil de verre mince
fait d'un \emph{cœur}\index{cœur (fibre optique)} d'indice $n_1$
enveloppé d'une \emph{gaine}\index{gaine (fibre optique)} d'indice
légèrement plus petit $n_2 < n_1$. La lumière injectée à une extrémité
voyage le long du cœur par \omterm{def:g10:refraction:critical}{réflexion totale interne} répétée sur
la frontière cœur--gaine.
\end{definition}

\begin{omfigure}
\begin{tikzpicture}[scale=1.0, font=\small]
  \draw[omExo, thick] (0,0.85) -- (7,0.85);
  \draw[omExo, thick] (0,-0.85) -- (7,-0.85);
  \fill[blue!8] (0,-0.5) rectangle (7,0.5);
  \draw[thick] (0,0.5) -- (7,0.5);
  \draw[thick] (0,-0.5) -- (7,-0.5);
  \draw[-Stealth, omThm, thick] (0,-0.36) -- (1.0,0.5) -- (2.16,-0.5)
    -- (3.32,0.5) -- (4.48,-0.5) -- (5.64,0.5) -- (6.6,-0.33);
  \node[anchor=south] at (3.5,0.9) {gaine ($n_2 = 1.46$)};
  \node[fill=blue!8, inner sep=2pt] at (3.5,-0.25) {cœur ($n_1 = 1.48$)};
\end{tikzpicture}

{\small Lumière guidée le long du \omterm{def:g10:refraction:fiber}{cœur} d'une fibre par
\omterm{def:g10:refraction:critical}{réflexion totale interne} (angles de rebond fortement exagérés~:
les vrais rayons rasent la paroi à plus de $80^\circ$ de la
\omterm{def:g10:refraction:angles}{normale}).}
\end{omfigure}

\begin{example}[Une fibre en chiffres]\label{ex:g10:refraction:fibernumbers}
Avec $n_1 = 1.48$ et $n_2 = 1.46$,
\[
\sin i_c = \frac{1.46}{1.48} = 0.986,
\qquad i_c = 80.6^\circ :
\]
seuls les rayons à environ $9^\circ$ de l'axe sont guidés --- ce qui
convient parfaitement, puisque c'est exactement ainsi que la lumière
laser est injectée. Dans le \omterm{def:g10:refraction:fiber}{cœur} la lumière voyage à
$v = c/1.48 = \qty{2.03e8}{m/s}$~: une impulsion traverse une fibre de
\qty{100}{km} en environ une demi-milliseconde.
\end{example}

\begin{remark}[Ordre de grandeur des liaisons de données]\label{rem:g10:refraction:datarates}
Les câbles sous-marins --- faisceaux de quelques dizaines de fibres sur
le fond de l'océan --- portent presque tout le trafic intercontinental.
Une seule fibre moderne transporte de l'ordre de $10^{13}$~bits par
seconde en envoyant de nombreuses \omterm{def:g10:light-spectra:wavelength}{longueurs d'onde} à la fois~; un
câble entier atteint $10^{14}$~bits par seconde, avec un amplificateur
tous les \qty{80}{km} environ. Les satellites, pour comparaison,
gèrent une petite fraction d'un pourcent du trafic.
\end{remark}

\section{Exercices}

\begin{exercise}[$\star$]\label{exo:g10:refraction:1}
Un faisceau laser frappe un miroir plan à $25^\circ$ \emph{depuis la
surface du miroir}. Donner l'\omterm{def:g10:refraction:angles}{angle d'incidence}, l'angle de
réflexion, et l'angle entre les faisceaux incident et réfléchi.
\end{exercise}

\begin{exercise}[$\star$]\label{exo:g10:refraction:2}
\begin{enumerate}
  \item Calculer la vitesse de la lumière dans l'eau ($n = 1.33$) et
        dans le diamant ($n = 2.42$).
  \item Dans un certain solide, la lumière voyage à \qty{1.97e8}{m/s}.
        Calculer son \omterm{def:g10:refraction:index}{indice de réfraction} et identifier le matériau.
\end{enumerate}
\end{exercise}

\begin{exercise}[$\star$]\label{exo:g10:refraction:3}
Un rayon dans l'air frappe un bloc de verre ($n = 1.50$) à
$i_1 = 40^\circ$. Calculer l'angle de \omterm{def:g10:refraction:refraction}{réfraction}. Sous quel angle un
rayon voyageant à l'intérieur du verre devrait-il frapper la surface
pour sortir dans l'air à $40^\circ$~?
\end{exercise}

\begin{exercise}[$\star$]\label{exo:g10:refraction:4}
La lumière du Soleil atteint la surface d'un lac à $55^\circ$ de la
verticale. Calculer l'angle du \omterm{def:g10:refraction:refraction}{rayon réfracté}, et indiquer si le
rayon se courbe vers ou en s'écartant de la \omterm{def:g10:refraction:angles}{normale} --- avec la
raison en un mot.
\end{exercise}

\begin{exercise}[$\star$]\label{exo:g10:refraction:5}
Pour identifier un liquide transparent, un rayon est envoyé sur sa
surface à $i_1 = 45^\circ$~; le \omterm{def:g10:refraction:refraction}{rayon réfracté} est mesuré à
$i_2 = 32^\circ$. Calculer l'\omterm{def:g10:refraction:index}{indice de réfraction} et identifier le
liquide d'après le tableau de la~\cref{def:g10:refraction:index}.
\end{exercise}

\begin{exercise}[$\star\star$]\label{exo:g10:refraction:6}
Un demi-disque d'un plastique transparent donne les mesures suivantes~:
\begin{center}
\begin{tabular}{lccccc}
$i_1$ & $20^\circ$ & $30^\circ$ & $40^\circ$ & $50^\circ$ & $60^\circ$ \\
$i_2$ & $13.5^\circ$ & $19.5^\circ$ & $25.5^\circ$ & $31.0^\circ$
      & $35.5^\circ$ \\
\end{tabular}
\end{center}
\begin{enumerate}
  \item Calculer $\dfrac{\sin i_1}{\sin i_2}$ pour chaque paire. Que
        constatez-vous~?
  \item En déduire l'\omterm{def:g10:refraction:index}{indice de réfraction} du plastique et l'identifier.
  \item Pourquoi tracer $\sin i_2$ en fonction de $\sin i_1$ est-il un
        meilleur test de la \omterm{thm:g10:refraction:snell}{loi de Snell--Descartes} que tracer
        $i_2$ en fonction de $i_1$~?
\end{enumerate}
\end{exercise}

\begin{exercise}[$\star\star$]\label{exo:g10:refraction:7}
Calculer l'\omterm{def:g10:refraction:critical}{angle critique} vers l'air pour l'eau, le verre
($n = 1.50$) et le diamant. Puis expliquer, avec la
\omterm{thm:g10:refraction:snell}{loi de Snell--Descartes}, pourquoi il n'y a pas d'\omterm{def:g10:refraction:critical}{angle
critique} pour la lumière allant \emph{de} l'air \emph{vers} l'eau.
\end{exercise}

\begin{exercise}[$\star\star$]\label{exo:g10:refraction:8}
\begin{enumerate}
  \item Une piscine a \qty{2.0}{m} de profondeur. Quelle profondeur un
        nageur debout au bord perçoit-il~?
  \item Un héron vise un poisson nageant à \qty{40}{cm} sous la
        surface. À quelle profondeur le poisson apparaît-il, et le
        héron doit-il frapper au-dessus, sur, ou en dessous de l'image
        qu'il voit~?
\end{enumerate}
\end{exercise}

\begin{exercise}[$\star\star$]\label{exo:g10:refraction:9}
Un rayon frappe une vitre ($n = 1.50$, faces parallèles) à $50^\circ$.
\begin{enumerate}
  \item Calculer l'angle de \omterm{def:g10:refraction:refraction}{réfraction} à la première face.
  \item Le \omterm{def:g10:refraction:refraction}{rayon réfracté} frappe la seconde face de l'intérieur, au
        même angle. Calculer l'angle de sortie et le comparer à
        $50^\circ$.
  \item Que fait alors une vitre épaisse à la direction et à la
        position d'un rayon qui la traverse~?
\end{enumerate}
\end{exercise}

\begin{exercise}[$\star\star$]\label{exo:g10:refraction:10}
De la \omterm{def:g10:light-spectra:white}{lumière blanche} frappe un bloc de flint glass à $55^\circ$.
Pour ce verre, $n_{\text{red}} = 1.61$ et $n_{\text{blue}} = 1.66$.
Calculer les angles de \omterm{def:g10:refraction:refraction}{réfraction} des rayons rouge et bleu et
l'angle entre eux. Quelle couleur finit plus près de la
\omterm{def:g10:refraction:angles}{normale}~?
\end{exercise}

\begin{exercise}[$\star\star$]\label{exo:g10:refraction:11}
Une \omterm{def:g10:refraction:fiber}{fibre optique} a un \omterm{def:g10:refraction:fiber}{cœur} d'indice $1.48$ et une
\omterm{def:g10:refraction:fiber}{gaine} d'indice $1.46$.
\begin{enumerate}
  \item Calculer l'\omterm{def:g10:refraction:critical}{angle critique} sur la frontière cœur--gaine.
  \item Un rayon guidé rebondit exactement à l'\omterm{def:g10:refraction:critical}{angle critique}.
        Montrer que sur \qty{1.0}{km} de fibre droite il parcourt environ
        \qty{14}{m} de plus qu'un rayon allant droit le long de l'axe.
\end{enumerate}
\end{exercise}

\begin{exercise}[$\star\star\star$]\label{exo:g10:refraction:12}
L'œil d'un plongeur est à \qty{2.0}{m} sous une surface parfaitement
calme. À cause de la \omterm{def:g10:refraction:refraction}{réfraction}, le plongeur voit le ciel
\emph{entier} comprimé en un disque brillant juste au-dessus --- la
fenêtre de Snell.
\begin{enumerate}
  \item Expliquer pourquoi la lumière de l'horizon atteint l'œil du
        plongeur à l'\omterm{def:g10:refraction:critical}{angle critique} de l'eau, $48.8^\circ$ de la
        verticale.
  \item Calculer le rayon du disque brillant sur la surface, en
        utilisant la tangente dans le triangle rectangle
        œil--verticale--bord du disque.
  \item Que voit le plongeur sur la surface \emph{à l'extérieur} du
        disque~?
\end{enumerate}
\end{exercise}

\begin{exercise}[$\star\star\star$]\label{exo:g10:refraction:13}
Un prisme a des angles $30^\circ$, $60^\circ$, $90^\circ$. Un faisceau
mince de \omterm{def:g10:light-spectra:white}{lumière blanche} entre perpendiculairement à l'un des côtés
adjacents à l'angle droit, traverse sans déviation, et frappe la grande
face de l'intérieur à $i_1 = 30^\circ$. Pour ce verre,
$n_{\text{red}} = 1.51$ et $n_{\text{blue}} = 1.53$.
\begin{enumerate}
  \item Calculer les angles de sortie des rayons rouge et bleu et la
        largeur angulaire de l'éventail entre eux.
  \item Le prisme est remplacé par un prisme $45^\circ$--$45^\circ$--$90^\circ$,
        de sorte que l'\omterm{def:g10:refraction:angles}{angle d'incidence} interne devient $45^\circ$.
        Montrer qu'aucune lumière ne sort par la grande face.
\end{enumerate}
\end{exercise}

\begin{exercise}[$\star\star\star$]\label{exo:g10:refraction:14}
Les périscopes plient la lumière avec des prismes de verre
$45^\circ$--$45^\circ$--$90^\circ$ ($n = 1.50$)~: le faisceau entre par
une petite face perpendiculairement et frappe la grande face de
l'intérieur à $45^\circ$.
\begin{enumerate}
  \item Montrer que la grande face agit comme un miroir parfait.
        Pourquoi est-ce mieux qu'un miroir métallique, qui absorbe
        quelques pourcents à chaque réflexion~?
  \item Le périscope se remplit d'eau et la grande face est maintenant
        en contact avec l'eau. Calculer le nouvel \omterm{def:g10:refraction:critical}{angle critique}
        de la frontière verre--eau et montrer que le miroir échoue.
  \item Sous quel angle la lumière de la question~2 quitte-t-elle le
        prisme vers l'eau~?
\end{enumerate}
\end{exercise}

\begin{exercise}[$\star\star\star$]\label{exo:g10:refraction:15}
Un câble transatlantique court sur \qty{6200}{km} de fibre (\omterm{def:g10:refraction:fiber}{cœur}
d'indice $1.47$) entre Londres et New York.
\begin{enumerate}
  \item Calculer la vitesse de la lumière dans le \omterm{def:g10:refraction:fiber}{cœur} et le temps
        de parcours aller d'une impulsion.
  \item Une onde radio voyage à $c$. Combien de temps lui faudrait-il
        sur la même distance, et combien de temps aller-retour le verre
        de la fibre coûte-t-il par rapport à la radio~?
  \item Certaines firmes financières paient des fortunes pour gratter
        des millisecondes sur ce lien. Proposer deux façons physiques
        distinctes de faire arriver la lumière plus tôt.
\end{enumerate}
\end{exercise}

\section{Problème~: Le diamant, la piscine et la fibre}

\begin{problem}[{Problème du week-end --- trois vies de la loi de
Snell--Descartes~: une piscine qui ment sur sa profondeur, une pierre
taillée pour que la lumière ne puisse en sortir, et l'océan de soixante
millisecondes sous l'internet}]
\label{pb:g10:refraction:1}
Une égalité de deux sinus, trois carrières. Dans une piscine elle
trompe vos yeux sur la profondeur et comprime tout le ciel en un
cercle. Dans un diamant, poussée à sa limite, elle refuse de laisser
sortir la lumière --- et les joailliers taillent les pierres autour de
ce refus depuis des siècles. Sous l'Atlantique elle guide des
impulsions laser dans un brin de verre, et fixe la limite de vitesse de
l'internet mondial. Ce problème parcourt les trois vies dans l'ordre,
avec la même loi à chaque fois.

\medskip
\noindent\textbf{Partie I --- Mise en train.}
\begin{enumerate}
  \item Calculer la vitesse de la lumière dans l'eau ($n = 1.33$) et
        dans le diamant ($n = 2.42$). D'un facteur combien le diamant
        ralentit-il la lumière~?
  \item Un rayon dans l'air frappe un étang à $i_1 = 40^\circ$.
        Calculer l'angle de \omterm{def:g10:refraction:refraction}{réfraction}.
  \item Un rayon dans l'eau frappe la surface d'en dessous à
        $28.9^\circ$. Sous quel angle sort-il dans l'air~? Énoncer le
        principe que cela illustre.
  \item Calculer l'\omterm{def:g10:refraction:critical}{angle critique} de la frontière eau--air, et
        décrire ce qui arrive à un rayon sous l'eau frappant la surface
        à $60^\circ$.
  \item L'indice de l'air est $1.0003$. Calculer la vitesse de la
        lumière dans l'air et justifier l'approximation
        $n_{\text{air}} = 1.00$ utilisée partout.
\end{enumerate}

\medskip
\noindent\textbf{Partie II --- La piscine.}
\begin{enumerate}[resume]
  \item Une piscine a \qty{3.0}{m} de profondeur. En utilisant la
        formule de profondeur apparente
        (\cref{prop:g10:refraction:depth}), calculer la profondeur qu'une
        personne debout au-dessus perçoit.
  \item Expliquer avec un croquis à deux rayons (décrit en mots)
        pourquoi une pièce au fond semble monter~: que font les rayons
        quittant l'eau, et où l'œil place-t-il leur intersection~?
  \item L'œil d'un plongeur est à \qty{3.0}{m}. Calculer le rayon de la
        fenêtre de Snell --- le disque brillant à travers lequel le
        plongeur voit tout le ciel (\cref{exo:g10:refraction:12}).
  \item Une lampe au fond envoie des rayons vers le haut à $30^\circ$,
        $45^\circ$ et $50^\circ$ de la verticale. Pour chaque rayon,
        dire s'il s'échappe et sous quel angle, ou ce qui se passe à la
        place.
  \item Un pêcheur à la lance sur la berge vise un poisson. Le coup
        doit-il être porté au-dessus, sur, ou en dessous de l'image du
        poisson --- et le poisson, regardant en arrière, voit-il aussi
        le pêcheur déplacé~?
\end{enumerate}

\medskip
\noindent\textbf{Partie III --- Le diamant.}
\begin{enumerate}[resume]
  \item Calculer l'\omterm{def:g10:refraction:critical}{angle critique} du diamant ($n = 2.42$) vers
        l'air, et le comparer à celui du verre ($n = 1.52$), qui est
        $41.1^\circ$.
  \item Un rayon à l'intérieur de la pierre frappe une facette à
        $30^\circ$ de sa \omterm{def:g10:refraction:angles}{normale}. Que se passe-t-il dans le diamant~?
        Dans l'imitation en verre~? Expliquer en une phrase pourquoi la
        taille d'un diamant, combinée à son minuscule \omterm{def:g10:refraction:critical}{angle
        critique}, produit l'éclat.
  \item Le diamant est aussi fortement \omterm{def:g10:refraction:dispersion}{dispersif}~:
        $n_{\text{red}} = 2.41$, $n_{\text{blue}} = 2.45$. De la
        \omterm{def:g10:light-spectra:white}{lumière blanche} entre dans une facette à $45^\circ$~;
        calculer l'angle de \omterm{def:g10:refraction:refraction}{réfraction} pour chaque couleur et
        l'écart angulaire.
  \item Plus tard, des rayons rouge et bleu atteignent tous deux une
        facette de sortie de l'intérieur à $17.0^\circ$. Calculer les
        deux angles de sortie et l'écart de l'éventail émergent.
        Comparer à l'écart à l'entrée~: que devient la
        \omterm{def:g10:refraction:dispersion}{dispersion} près de l'\omterm{def:g10:refraction:critical}{angle critique}~?
  \item Une goutte d'eau ($n = 1.33$) enduit une facette. Calculer le
        nouvel \omterm{def:g10:refraction:critical}{angle critique} de la frontière diamant--eau, décider
        du sort du rayon à $30^\circ$ de la question~12, et expliquer
        pourquoi les joailliers gardent les diamants scrupuleusement
        propres.
\end{enumerate}

\medskip
\noindent\textbf{Partie IV --- La fibre.}
\begin{enumerate}[resume]
  \item Une fibre a un \omterm{def:g10:refraction:fiber}{cœur} d'indice $1.48$ et une
        \omterm{def:g10:refraction:fiber}{gaine} d'indice $1.46$. Calculer l'\omterm{def:g10:refraction:critical}{angle critique} sur
        la frontière cœur--gaine.
  \item Un fil de verre nu dans l'air aurait un \omterm{def:g10:refraction:critical}{angle critique} de
        $42.5^\circ$ --- apparemment bien meilleur. Donner deux raisons
        pour lesquelles la \omterm{def:g10:refraction:fiber}{gaine} est utilisée quand même. (Penser à
        la poussière, aux rayures, et à ce que le toucher de la surface
        ferait à la \omterm{def:g10:refraction:critical}{réflexion totale interne}.)
  \item Calculer la vitesse de la lumière dans le \omterm{def:g10:refraction:fiber}{cœur}, puis le
        temps de parcours aller d'une impulsion le long de
        \qty{6000}{km} de fibre.
  \item Un rayon rebondissant exactement à l'\omterm{def:g10:refraction:critical}{angle critique}
        parcourt $1/\sin i_c$ fois la longueur de l'axe. Calculer la
        distance supplémentaire sur les \qty{6000}{km} et le temps
        supplémentaire qu'elle coûte.
  \item La chute. Une paire de fibres moderne porte environ $10^{13}$
        bits par seconde, et ce livre fait environ
        $2.4 \times 10^8$~bits. Calculer le temps aller-retour
        («~ping~») à travers l'océan et le nombre de copies de ce livre
        que la fibre déplace par seconde --- puis énoncer, en une
        phrase, ce que la \omterm{thm:g10:refraction:snell}{loi de Snell--Descartes} fait pour
        l'internet.
\end{enumerate}
\end{problem}
